\documentclass[11pt,a4paper]{article}
\usepackage[margin=1in]{geometry}
\usepackage[T1]{fontenc}
\usepackage{lmodern,amsmath,amssymb,amsthm,booktabs,array,microtype}
\usepackage[colorlinks=true,linkcolor=blue,urlcolor=blue]{hyperref}
\hypersetup{pdftitle={Verification of the Prophet Secretary Certificate},
 pdfsubject={A finite certificate for the 0.6916 mean-only Prophet Secretary guarantee}}
\newtheorem{theorem}{Theorem}
\newtheorem{lemma}{Lemma}
\newtheorem{proposition}{Proposition}
\newcommand{\ind}{\mathbf{1}}
\newcommand{\pos}[1]{\left(#1\right)_+}
\newcommand{\Hp}{H_{p^*}}
\newcommand{\Sh}{\widehat S}
\setlength{\emergencystretch}{2em}
\title{Verification of the Prophet Secretary Certificate}
\author{}
\date{}
\begin{document}
\maketitle
\begin{abstract}
This supplement gives the coefficient construction and the numerical enclosure
arguments for the finite certificate accompanying \emph{Mean-only Prophet Inequalities}. It uses the manuscript's original
$\Gamma_{p^*}$, Lemmas 4.7 and 4.8, and the discretization in Lemmas B.1 and B.2.
The supplied rational threshold and proof parameters certify a uniform lower
bound of $49839735470020217/72057594037927936>0.6916$.
We specify the bounds on complete time intervals, the treatment of actual
probability averages, and the directed arithmetic used by the verifier.
The source of this document is one self-contained \texttt{Verification.tex}
file; no external mathematical source files are needed to compile it.
\end{abstract}

\section{Statement and relation to the manuscript}
References to numbered lemmas in the manuscript below refer to the version in
which the Prophet Secretary analysis is Section 4 and its finite LP is
Appendix B.2. We use the same notation $b,\theta,p^*,P^*,H_{p^*},\widehat S$
and $\Gamma_{p^*}$. In particular,
\[
 \theta(t)=\int_t^1 b(s)\,ds,\qquad
 P^*(t)=\int_0^t p^*(v)\,dv,
 \qquad
 \widehat S(s)=(1-p^*(s))
       \exp\!\left(\int_0^s\log(1-p^*(v))\,dv\right).
\]
Here $p^*:[0,1]\to[0,1]$ is any nondecreasing function; it need not be
constant on a grid interval or continuous. Integrals are Lebesgue integrals.
Values at finitely many partition endpoints do not affect them. As in the
manuscript, logarithms at zero and the associated expressions are interpreted
by limits.

\begin{theorem}\label{thm:certificate}
The parameters in \texttt{parameters.json}, together with the fixed coefficient
choices in \texttt{proposals/}, define an admissible instance of the manuscript's
Prophet Secretary construction with
\[
 N=160,\qquad N_t=5120,\qquad N_p=10240.
\]
For these common parameters, every nondecreasing $p^*:[0,1]\to[0,1]$ and
every $\tau\in[0,1]$ satisfy
\begin{equation}\label{eq:certified}
 \Gamma_{p^*}(\tau)\ \ge\
 \frac{49839735470020217}{72057594037927936}
 =0.6916652732505443\ldots\ >\ \frac{1729}{2500}=0.6916.
\end{equation}
The threshold is continuous, concave and strictly decreasing, with
\[
 \theta(0)=\frac{567246577807}{800000000000}
          =0.70905822225875.
\]
Consequently, the manuscript's Lemma 4.6 yields the claimed $0.6916$
competitive ratio in Theorem 1.2 for every $n$.
\end{theorem}

The finite arithmetic checks used in this theorem are specified below and can
be repeated with the included verifier. The certificate establishes feasibility,
not optimality of the threshold or of the finite LP. It concerns the Prophet
Secretary lower bound; it is not a certificate for the manuscript's separate
i.i.d. or upper-bound results.

\section{Common data, grids and fixed plane choices}\label{sec:data}
The manuscript indexes drop rates by $1,\ldots,N$. For zero-based file index
$a\in\{0,\ldots,N-1\}$, put $J_a=[a/N,(a+1)/N)$.
The functions $b,m,\kappa$ are constant on $J_a$, and
\[
 \begin{aligned}
 b(s)=b_{a+1}&=\frac{\texttt{b\_numerators}[a]}{10^{10}},\\
 m(s)&=\frac{\texttt{m\_numerators}[a]}{10^{10}},\qquad
 \kappa(s)=\frac{\texttt{kappa\_numerators}[a]}{10^{10}}.
 \end{aligned}
\]
The auxiliary function $m$ specifies the allocation between the two original
extra-gain bounds; its precise use is given in Section~\ref{sec:allocation}.
The verifier checks
\[
 0<b_1<\cdots<b_N,\qquad
 \frac1N\sum_{a=1}^N b_a\le1,\qquad
 0\le m(s)\le b(s),\qquad \kappa(s)\ge0.
\]
Every proposal contains the same 640 common integer parameters. The first
three blocks specify these three functions; the fourth block is checked to
be identically zero and contributes no term.

Write
\[
 h=\frac1{N_p},\qquad s_j=jh,\qquad I_j=[s_j,s_{j+1}],\qquad
 \xi_j=\frac1h\int_{I_j}p^*(s)\,ds\quad(0\le j<N_p).
\]
Monotonicity gives $0\le\xi_0\le\cdots\le\xi_{N_p-1}\le1$.
The time intervals are
\[
 T_\ell=[t_\ell^-,t_\ell^+]
       =[2\ell h,(2\ell+2)h],\qquad 0\le\ell<N_t.
\]
They cover $[0,1]$. Each lies within a single constant-slope piece of
$\theta$. The probability partition refines all breakpoints of $b,m,\kappa$
and the route functions used below.

Each proposal has an additional plane-owner partition into equal cells.
Its size is 160, 320 or 640 in this release, and divides $N_p$.
It specifies three types of plane lists: ordinary $H$ planes, capped $H$
planes, and $\widehat S$ planes. A plane has a nonnegative integer mass with
denominator $2^{32}$. The total mass is checked to be at most one separately
for every owner and every type. Those masses and the plane slopes or $\rho$
values are fixed on the entire $T_\ell$.
The two stored $H$ slope integers represent $\sigma$ and $\nu$ with
denominator $2^{32}$; a stored survival-plane parameter represents
$\rho$ with denominator $10^{10}$.

If $f\ge0$ and $f\ge L_i$ pointwise, then, for $\omega_i\ge0$ with
$\sum_i\omega_i\le1$,
\[
 f\ge\sum_i\omega_i L_i.
\]
Indeed, give unused mass to the zero lower bound. This applies to both
$H_{p^*}$ and $\widehat S$, which are nonnegative. For the former, when its
logarithms are finite,
\[
 \mathcal L(d,x,y)
 =-\log\!\left(1-\frac{dx}{1-y}\right)-dx\ge0
\]
by $-\log(1-z)\ge z$ for $0\le z<1$; the limiting cases follow as well.
Thus the plane mixtures are fixed affine lower bounds of exactly the kind
permitted by Lemma B.1. They do not mix algorithms or thresholds.

\paragraph{Deterministic assembly.}
The supplied generator chooses, for each complete $T_\ell$, a rational
$\alpha_\ell\in[0,1]$ from frozen input data. A source mass $M$ is replaced
by $\lfloor(1-\alpha_\ell)M\rfloor$ or $\lfloor\alpha_\ell M\rfloor$.
After refinement to a common owner partition, each child receives the
corresponding parent list. Since each source list has mass at most one,
the new total remains at most one. The verifier also checks this directly.
The choice may be guided by the midpoint of $T_\ell$, but the resulting
mixture is held fixed throughout that interval. Every $H$ intercept is then
recomputed for that whole rectangle. No interpolation of certified bounds
or time-dependent mixture is used.


\section{Allocation and comparison contributions}\label{sec:allocation}
\subsection{Fixed allocation and admissibility}
The endpoint and starting-point functions are fixed independently of
$p^*$.  For $s\in J_a$, define
\[
 u(s)=r_a+(1-r_a)s,\qquad z(s)=q_a s,
\]
where the following table specifies every route used by the verifier.
The first column is $\lfloor a/4\rfloor$, rather than the threshold-segment
index $a$ itself.
\begin{center}
\begin{tabular}{c|c|c}
$\lfloor a/4\rfloor$ & $r_a$ & $q_a$\\ \hline
$16$ & $3/4$ & $1$\\
$17$ & $7/10$ & $9/10$\\
$18,\ldots,27$ & $1/2$ & $1/2$\\
$32$ & $1/4$ & $1/2$\\
$33,\ldots,37$ & $1/4$ & $1/4$\\
all remaining indices & $0$ & $1$
\end{tabular}
\end{center}
In particular, $s\leq u(s)\leq1$ and $0\leq z(s)\leq s$.

In the notation of the article, use two applications of Lemma~4.3 and
one comparison application of Lemma~4.4:
\[
 \begin{aligned}
 c_1(s)&=1,&\lambda_1(s)&=b(s)-m(s),\\
 c_2(s)&=u(s),&\lambda_2(s)&=m(s),\\
 z_1(s,t)&=z(s),&\mu_1(s,t)&=m(s)\,\mathbf{1}_{\{t>u(s)\}}.
 \end{aligned}
\]
The strict inequality in the last indicator ensures that the allocation
constraint holds also at $t=u(s)$, where the capped application of
Lemma~4.3 still uses its coefficient.  For every $s\leq t<1$, the total
allocated coefficient is exactly $b(s)$: it is
$(b(s)-m(s))+m(s)$ before and at $u(s)$, and
$(b(s)-m(s))+m(s)$ after $u(s)$, with the second summand coming from
the comparison application in the latter case.  Hence (4.1) holds
pointwise.  Changing the value of an indicator at $t=u(s)$ leaves all
integrals below unchanged.

For these data, $m(s)$ is positive precisely on the union of the
threshold segments with indices $64\leq a\leq111$ and
$136\leq a\leq151$.  The route table consequently gives
\[
 m(s)>0\quad\Longrightarrow\quad z(s)\geq\frac{17}{80}.
\]
In particular, the comparison denominator satisfies
$1-t+z(s)\geq17/80$ wherever the comparison weight is nonzero.
Since $m$ is bounded, the comparison integrability condition in
Appendix~B.2 holds.  All changes in the order of integration below are
therefore justified; for the nonnegative terms they also follow directly
from Tonelli's theorem.  Expressions multiplied by $m(s)=0$ are taken
to contribute zero, so no value at an irrelevant zero denominator is
needed.

\subsection{Exact comparison kernels and conservative cell coefficients}
\label{cert:comparison}

Write $P^*(v)=\int_0^v p^*(w)\,dw$.  With the allocation just specified,
the comparison contribution to $\Gamma_{p^*}(\tau)$ is
\begin{equation}
 \int_0^1m(s)\int_{u(s)}^1
 \left(
   \mathbf{1}_{\{\tau\leq t\}}p^*(s)
   +\frac{P^*(t)-P^*(z(s))-\mathbf{1}_{\{\tau\leq t\}}}
          {1-t+z(s)}
 \right)\,dt\,ds.
 \label{cert:comparison-original}
\end{equation}
The equality uses only the preceding allocation and the article's
definition of $\Gamma_{p^*}$.
Together with
\[
 F_{p^*}(\tau)=\tau+\theta(\tau)
  -\int_0^1b(s)\min\{s,\tau\}p^*(s)\,ds,
\]
the complete instantiated criterion consists of $F_{p^*}(\tau)$,
the comparison contribution \eqref{cert:comparison-original}, and
\begin{align*}
 &\int_0^1(b(s)-m(s))H_{p^*}(s,(\tau-s)_+)\,ds\\
 &\quad+\int_0^1m(s)H_{p^*}(s,(\min\{\tau,u(s)\}-s)_+)\,ds\\
 &\quad+\int_0^1\kappa(s)\widehat S(s)\,ds
          -\int_0^\tau\kappa(s)(1-p^*(s))\,ds.
\end{align*}
Thus every term of the original criterion is accounted for below.

For $0\leq v\leq1$, define the nonnegative kernel
\begin{equation}
 K(v)=\int_0^1m(s)
 \log\!\left(\frac{1-\max\{v,u(s)\}+z(s)}{z(s)}\right)\,ds.
 \label{cert:kernel}
\end{equation}
Indeed, the logarithm equals
$\int_{\max\{v,u(s)\}}^1(1-t+z(s))^{-1}\,dt$.
It is nonincreasing in $v$ for each fixed source point $s$.

Fix one stopping-time interval
$T_\ell=[t_\ell^-,t_\ell^+]$ and a probability cell
$I_j=[s_j,s_{j+1}]$, of length $h=1/N_p$ with $N_p=10240$.
The common probability partition refines the threshold and route
partitions.  On $I_j$, use the affine extension of its own route and put
\[
 \overline u_j=r_a+(1-r_a)s_{j+1}
 \quad\text{when the interior of $I_j$ lies in $J_a$}.
\]
This definition uses the route on the interior of $I_j$, even if the
next cell uses a different route.  Thus $u(s)\leq\overline u_j$ for
almost every $s\in I_j$.
As in the article, write
\[
 \xi_j=\frac1h\int_{I_j}p^*(s)\,ds.
\]
The first term in \eqref{cert:comparison-original} integrates exactly to
$\int_0^1m(s)(1-\max\{\tau,u(s)\})p^*(s)\,ds$.
For every $\tau\in T_\ell$, its contribution from $I_j$ is at least
\begin{equation}
 h\,m(s_j)\bigl(1-\max\{t_\ell^+,\overline u_j\}\bigr)\xi_j.
 \label{cert:comparison-direct}
\end{equation}
Here and below the step functions are represented by their values on
the cell interior, with the right-continuous convention at $s_j$.
The term with the negative indicator is exactly $-K(\tau)$.
Since $K$ is nonincreasing, it is at least $-K(t_\ell^-)$ throughout
$T_\ell$.

For the remaining, nonnegative comparison-history term, use
$P^*(t)-P^*(z(s))=\int_{z(s)}^t p^*(v)\,dv$ and exchange the integrals.
The result is
\begin{equation}
 \int_0^1p^*(v)
 \left[\int_{\{s:z(s)\leq v\}}m(s)
 \log\!\left(\frac{1-\max\{v,u(s)\}+z(s)}{z(s)}\right)\,ds\right]dv.
 \label{cert:comparison-fubini}
\end{equation}
For $v\in I_j$, retain only the source points with $z(s)\leq s_j$.
They are eligible for every $v$ in that cell.  On this fixed source set,
replace $v$ by $s_{j+1}$ inside the nonincreasing logarithmic kernel.
Define
\begin{equation}
 K_j=\int_{\{s:z(s)\leq s_j\}}m(s)
 \log\!\left(\frac{1-\max\{s_{j+1},u(s)\}+z(s)}{z(s)}\right)\,ds.
 \label{cert:retained-kernel}
\end{equation}
Then \eqref{cert:comparison-fubini} is at least
$h\sum_jK_j\xi_j$.  This argument does not assume that the expression
in square brackets in \eqref{cert:comparison-fubini} is nonincreasing:
its source set grows with $v$.  The source set in
\eqref{cert:retained-kernel} is fixed before taking a cell endpoint.

Combining the three contributions, a lower bound on
\eqref{cert:comparison-original}, uniform over $\tau\in T_\ell$, is
\begin{equation}
 -K(t_\ell^-)
 +h\sum_{j=0}^{N_p-1}
 \left[m(s_j)\bigl(1-\max\{t_\ell^+,\overline u_j\}\bigr)+K_j\right]\xi_j.
 \label{cert:comparison-cell-bound}
\end{equation}
Every coefficient is fixed on the whole stopping-time interval and is
independent of $p^*$.

\paragraph{Evaluation of the logarithmic kernels.}
On a source segment with route constants $r=r_a$ and $q=q_a$, put
$d=(v-r)/(1-r)$.  In the portion with $s\leq d$ the integrand in
\eqref{cert:kernel}, without its factor $m(s)$, is
\[
 \log\frac{1-v+qs}{qs};
\]
in the portion with $s\geq d$ it is
\[
 \log\frac{1-r+(q+r-1)s}{qs}.
\]
The splitting point is clipped to the source integration interval.
For \eqref{cert:retained-kernel}, that interval is additionally cut off
at $s=s_j/q$.  If $v=1$, both kernel integrals vanish directly.
For $v<1$, an antiderivative of
$\log((A+Cs)/(qs))$ is
\begin{equation}
 \begin{cases}
 \displaystyle
 \frac{(A+Cs)(\log(A+Cs)-1)}{C}
       -s(\log(qs)-1),& C\ne0,\\[6pt]
 \displaystyle s\log A-s(\log(qs)-1),& C=0.
 \end{cases}
 \label{cert:kernel-antiderivative}
\end{equation}
Differentiating proves the formula wherever the logarithmic arguments
are positive.  The two branches use $(A,C)=(1-v,q)$ and
$(A,C)=(1-r,q+r-1)$, respectively.  Endpoint limits use
$x\log x\to0$ as $x\downarrow0$; in particular the antiderivative at
$s=0$ is $A(\log A-1)/C$ when $A>0$ and $C\ne0$, and is zero in
the remaining applicable limiting cases.  These are exactly the
branches evaluated by \texttt{primitive} and \texttt{integrate}.


\section{The logarithmic contributions on complete rectangles}
\label{sec:release-h-rectangles}

Let $p^*:[0,1]\to[0,1]$ be nondecreasing, and use the paper's
definitions of $P^*$, $H_{p^*}$, and $\mathcal L$.  In particular,
\[
 P^*(s)=\int_0^s p^*(v)\,dv,
 \qquad
 \mathcal L(\Delta,x,y)=\log(1-y)-\log(1-\Delta x-y)-\Delta x.
\]
For $0\le\Delta\le1-s$, the manuscript defines
\[
 H_{p^*}(s,\Delta)
 =\mathcal L\!\left(\Delta,p^*(s),
        \int_{s-\delta}^{s}p^*(v)\,dv\right),
 \qquad \delta=\min\{s,(1-\Delta)/2\}.
\]
The case $\Delta=0$ is interpreted as $H_{p^*}(s,0)=0$.
Fix a probability cell $I_j=[s_j,s_{j+1}]$, where
$s_i=ih$ and $h=1/N_p$, and a complete stopping-time interval
$T_\ell=[t_\ell^-,t_\ell^+]$.  All bounds in this section hold for
every $s\in I_j$ and every $\tau\in T_\ell$.  At a route breakpoint,
endpoint values of the route mean the extensions of its formula on
the cell under consideration; the choices at individual breakpoints
do not affect any integral.

\subsection{Duration and retained history}

For the ordinary logarithmic contribution, the duration is
$(\tau-s)_+$.  For the capped contribution, it is
$(\min\{\tau,u(s)\}-s)_+$, where $s\le u(s)\le1$ and the fixed
route $u$ is nondecreasing on the cell.  For the contribution being
considered, denote its duration by $\Delta(s,\tau)$ and set
\[
\begin{array}{c|c|c}
 & d & d_{\mathrm{up}}\\ \hline
\text{ordinary}
 &(t_\ell^- -s_{j+1})_+
 &(t_\ell^+ -s_j)_+\\[2pt]
\text{capped}
 &(\min\{t_\ell^-,u(s_j)\}-s_{j+1})_+
 &(\min\{t_\ell^+,u(s_{j+1})\}-s_j)_+
\end{array}
\]
Monotonicity of the positive-part and minimum operations gives
\[
 0\le d\le\Delta(s,\tau)\le d_{\mathrm{up}}\le1.
\]
If $d=0$, the verifier uses the zero lower bound for this
contribution.  This is valid because, wherever the logarithms are
finite,
\[
 \mathcal L(\Delta,x,y)
 =-\log\left(1-\frac{\Delta x}{1-y}\right)-\Delta x
 \ge \frac{\Delta x}{1-y}-\Delta x\ge0,
\]
and the same conclusion holds at the extended-real boundary.
We therefore consider $d>0$ below.

In the definition of $H_{p^*}(s,\Delta(s,\tau))$, the history length is
\[
 \delta(s,\tau)=
 \min\left\{s,\frac{1-\Delta(s,\tau)}2\right\}.
\]
Define the common lower bound
\[
 \delta_0=\min\left\{s_j,\frac{1-d_{\mathrm{up}}}2\right\},
 \qquad
 k=\min\left\{j,\max\left\{0,
   \left\lceil N_p(s_{j+1}-\delta_0)\right\rceil
 \right\}\right\}.
\]
The verifier retains the cells with indices $k,\ldots,j-1$, and puts
\[
 w=(j-k)h,
 \qquad
 y=\int_{s_k}^{s_j}p^*(v)\,dv
   =h\sum_{i=k}^{j-1}\xi_i,
 \qquad
 \xi_i=\frac1h\int_{I_i}p^*(s)\,ds.
\]
When $k=j$, the retained history is empty and both $w$ and $y$ are
zero; the displayed integral and sum are understood to be zero.
In particular, clipping a proposed starting index larger than $j$
does not retain an earlier nonempty interval.

If $k<j$, the unclipped ceiling is at most $j-1$, and
\[
 s_k\ge s_{j+1}-\delta_0,
 \qquad
 w=s_j-s_k\le\delta_0.
\]
Furthermore, $\delta(s,\tau)\ge\delta_0$.  Hence
\[
 s-\delta(s,\tau)
 \le s_{j+1}-\delta_0
 \le s_k\le s_j\le s.
\]
Thus the retained history lies inside the history in the definition
of $H_{p^*}$, simultaneously throughout the rectangle.  In the empty
case the same integral inequality holds directly.  We obtain
\[
 0\le y\le
 \int_{s-\delta(s,\tau)}^s p^*(v)\,dv,
 \qquad
 0\le y\le w\,p^*(s).
\]
The second inequality uses the monotonicity of the original function
$p^*$, since every retained point precedes $s$.

The logarithmic expression is nondecreasing in its duration and
history arguments.  Indeed, on its finite domain,
\[
 \frac{\partial\mathcal L}{\partial\Delta}
 =\frac{x}{1-\Delta x-y}-x\ge0,
 \qquad
 \frac{\partial\mathcal L}{\partial y}
 =\frac1{1-\Delta x-y}-\frac1{1-y}\ge0.
\]
Consequently,
\[
 H_{p^*}(s,\Delta(s,\tau))
 \ge \mathcal L(d,p^*(s),y).
\]
This inequality also holds when an original logarithm is singular,
by its limiting interpretation; such a value is $+\infty$ and causes
no difficulty for a finite lower bound.

The reduced triangle satisfies the strict domain requirement of
Lemma~4.7.  If $d_{\mathrm{up}}<1$, then
\[
 d+w\le d_{\mathrm{up}}+\frac{1-d_{\mathrm{up}}}2
       =\frac{1+d_{\mathrm{up}}}2<1.
\]
If $d_{\mathrm{up}}=1$, then $\delta_0=w=0$, while
$d\le(t_\ell^- -s_{j+1})_+<1$.  Thus $d+w<1$ also in this case.
Lemma~4.7 therefore applies with $\Delta=d$, $\delta=w$, and
\[
 D=\{(x,y):0\le x\le1,\ 0\le y\le wx\}.
\]

\subsection{Recomputing the intercept by the paper's edge tangents}

Fix the proposed slopes $\sigma,\nu\in\mathbb R$.  The verifier
recomputes an intercept for the reduced triangle, rather than
reusing an intercept from another rectangle.  For completeness,
write
\[
 R(x,y)=\mathcal L(d,x,y)-\sigma x-\nu y.
\]
If $w>0$, then $R$ is continuous on $D$ and
\[
 \det\nabla^2R(x,y)
 =-\frac{d^2}{(1-y)^2(1-dx-y)^2}<0
\]
at each interior point.  An interior local minimum would have a
positive semidefinite Hessian, so a global minimum occurs on the
boundary.  Parametrize the three edges by
\[
 (x_0(v),y_0(v))=(v,0),\qquad
 (x_1(v),y_1(v))=(v,wv),\qquad
 (x_2(v),y_2(v))=(1,wv),\quad 0\le v\le1,
\]
and put $\phi_e(v)=R(x_e(v),y_e(v))$.  As in the paper,
\[
\begin{aligned}
 \phi_0''(v)&=\frac{d^2}{(1-dv)^2}\ge0,\\
 \phi_1''(v)&=\frac{(d+w)^2}{(1-(d+w)v)^2}
              -\frac{w^2}{(1-wv)^2}\ge0,\\
 \phi_2''(v)&=\frac{w^2}{(1-d-wv)^2}
              -\frac{w^2}{(1-wv)^2}\ge0.
\end{aligned}
\]
For each edge, choose any $v_e\in[0,1]$.  Convexity gives the
uniform tangent lower bound
\[
 \phi_e(v)\ge\phi_e(v_e)+\phi_e'(v_e)(v-v_e),
 \qquad 0\le v\le1.
\]
It follows that
\[
 \beta_e=\phi_e(v_e)-v_e\phi_e'(v_e)
              +\min\{0,\phi_e'(v_e)\},
 \qquad
 \beta=\min_{e\in\{0,1,2\}}\beta_e
\]
is a valid intercept.  The selected tangent point need not minimize
$\phi_e$; its only required property is $v_e\in[0,1]$.
If $w=0$, the triangle is the segment $(v,0)$ and edge~0 suffices.
Including the duplicated or collapsed edges can only replace the
intercept by another valid lower bound.

We have proved, for the fixed rectangle and plane,
\[
 H_{p^*}(s,\Delta(s,\tau))
 \ge\beta+\sigma p^*(s)
          +\nu h\sum_{i=k}^{j-1}\xi_i.
\]
Averaging over the probability cell gives
\[
 \frac1h\int_{I_j}H_{p^*}(s,\Delta(s,\tau))\,ds
 \ge\beta+\sigma\xi_j
          +\nu h\sum_{i=k}^{j-1}\xi_i,
 \qquad \tau\in T_\ell.
\]
All averages here are averages of the actual $p^*$.  No step-function
replacement of $p^*$ has been made.  Multiplication by the fixed
nonnegative allocation on the cell and by any fixed nonnegative
plane mass preserves this inequality.

\section{The survival contribution and its coefficient error}
\label{sec:release-survival}

Fix $I_j=[s_j,s_{j+1}]$ and a plane parameter $\rho\in[0,1]$.
Apply Lemma~4.8 with $c=s_{j+1}>0$, and abbreviate
\[
 A=\rho^{s_{j+1}}(1+s_{j+1}-s_{j+1}\rho),
 \qquad B=\rho^{1+s_{j+1}}.
\]
With the paper's notation $\widehat S(s)$, the resulting bound is
\[
 \widehat S(s)\ge A(1-p^*(s))-B P^*(s),
 \qquad s\in I_j.
\]
Both constants are nonnegative.  When $\rho=0$, they are zero and
the following formulas remain valid.

Fubini's theorem on the bounded cell gives the exact identity
\[
 \frac1h\int_{s_j}^{s_{j+1}}P^*(s)\,ds
 =h\sum_{i<j}\xi_i
   +\frac1h\int_{s_j}^{s_{j+1}}(s_{j+1}-v)p^*(v)\,dv.
\]
For the last term, use the coefficient-error inequality in the
proof of Lemma~B.1 with coefficient
\[
 k(v)=-\frac{B}{h}(s_{j+1}-v),
 \qquad d_j=-\frac B2.
\]
Specifically, $0\le p^*\le1$ implies
\[
 \int_{I_j}k(v)p^*(v)\,dv
 \ge h d_j\xi_j-\int_{I_j}(d_j-k(v))_+\,dv.
\]
The error is exactly
\[
\begin{aligned}
 \int_{I_j}(d_j-k(v))_+\,dv
 &=\frac{B}{h}
   \int_{s_j}^{s_j+h/2}(s_j+h/2-v)\,dv\\
 &=\frac{Bh}{8}.
\end{aligned}
\]
Combining these expressions proves the cell-average lower bound
\[
 \boxed{
 \frac1h\int_{I_j}\widehat S(s)\,ds
 \ge A-\frac{Bh}{8}
       -\left(A+\frac{Bh}{2}\right)\xi_j
       -Bh\sum_{i<j}\xi_i.}
\]
Thus a fixed nonnegative correction allocation and plane mass add
their product with $h(A-Bh/8)$ to the constant, subtract their product
with $h(A+Bh/2)$ from the current-cell coefficient, and subtract
their product with $Bh^2$ from every preceding-cell coefficient.
This uses exactly the coefficient-error inequality stated in
Lemma~B.1, including its $Bh/8$ deduction before multiplication by
the cell length and allocation.

\section{Collecting coefficients and verifying the original LP}\label{sec:suffix}
For a fixed time interval $T_\ell$, the preceding bounds give exact real
coefficients $C_\ell,g_{\ell,0},\ldots,g_{\ell,N_p-1}$ such that
\begin{equation}\label{eq:affine}
 \Gamma_{p^*}(\tau)\ge C_\ell+\sum_{j=0}^{N_p-1}g_{\ell,j}\xi_j
 \qquad(\tau\in T_\ell).
\end{equation}
All plane parameters and route functions are fixed. These coefficients are
linear in the common drop rates and proof-weight coefficients, with the
baseline variable $v_\ell$ as in Lemma B.1. The common rational input is
one candidate solution of that original finite LP.

Here are the remaining baseline and prefix bounds used in collecting
\eqref{eq:affine}. Since $\tau+\theta(\tau)$ is affine on $T_\ell$, its
minimum is
\[
 v_\ell=\min\{t_\ell^-+\theta(t_\ell^-),
                  t_\ell^++\theta(t_\ell^+)\}.
\]
For the negative part of $F_{p^*}$, $b$ is constant on $I_j$, so
\[
 -\int_{I_j}b(s)\min\{s,\tau\}p^*(s)\,ds
 \ge -h b(s_j)\min\{s_{j+1},t_\ell^+\}\xi_j.
\]
For the correction prefix, use $\kappa(1-p^*)\ge0$ to obtain
\begin{align*}
 -\int_0^\tau\kappa(s)(1-p^*(s))\,ds
 &\ge-\int_0^{t_\ell^+}\kappa(s)(1-p^*(s))\,ds\\
 &=-\int_0^{t_\ell^+}\kappa(s)\,ds
        +\sum_{j:\ s_{j+1}\le t_\ell^+}h\kappa(s_j)\xi_j.
\end{align*}
The last equality uses the aligned grid, so no within-cell approximation
is present.

For completeness, Lemma B.2 eliminates the ordered averages exactly.
Set $\xi_{-1}=0$ and $\xi_{N_p}=1$. Then
\[
 C_\ell+\sum_j g_{\ell,j}\xi_j
 =\sum_{k=0}^{N_p}(\xi_k-\xi_{k-1})
       \left(C_\ell+\sum_{j=k}^{N_p-1}g_{\ell,j}\right).
\]
The coefficients form a convex combination. Conversely each suffix value
is attained by setting $\xi_j=\ind_{\{j\ge k\}}$. Hence the minimum is
\begin{equation}\label{eq:suffix}
 C_\ell+\min_{0\le k\le N_p}\sum_{j=k}^{N_p-1}g_{\ell,j}.
\end{equation}
The empty suffix $k=N_p$ is included. The verifier evaluates every suffix
in every time interval, a total of $5120\cdot10241=52\,433\,920$
suffix constraints.

The directed arithmetic described next supplies lower endpoints for
$C_\ell$ and every $g_{\ell,j}$. Since $\xi_j\ge0$, substituting those lower
endpoints into \eqref{eq:affine} remains a lower bound. Thus the numerical
check certifies all exact original LP constraints at the supplied rational
candidate; it does not require treating floating-point solver output or
rounded coefficients as a different mathematical LP.


\section{Directed arithmetic and certification of the numerical bounds}
\label{sec:directed-arithmetic}

All input parameters used in the proof are rational.  Set
\[
 Q=2^{56},\qquad G=2^{62}.
\]
The verifier represents an interval by two integers $(a,b)$, with
$-G<a\le b<G$, denoting $[a/Q,b/Q]$.  A rational number is enclosed by
rounding its scaled value down and up.  Addition and subtraction use the
corresponding endpoint formulas.  Multiplication takes the minimum and
maximum of the four endpoint products and divides by $Q$, rounding outwards.
Division similarly considers the four endpoint quotients; it is used only
when the denominator interval has strictly positive lower endpoint.
For a signed numerator and a positive denominator, the integer routines
compute the mathematical floor and ceiling, including when the numerator is
negative.  Consequently, each operation encloses the result for every choice
of real operands from its input intervals.  Shared occurrences of the same
quantity can enlarge the resulting interval, but cannot invalidate this
containment statement.

\paragraph{Logarithms.}
For a positive dyadic endpoint $x$, exact powers of two reduce $x$ to
$x=2^k r$, where $1\le r\le2$.  The verifier evaluates
$\log x=k\log2+\log r$ with directed arithmetic.  Writing
$z=(r-1)/(r+1)$ gives $0\le z\le1/3$ and
\[
 \log r=2\sum_{q=0}^{23}\frac{z^{2q+1}}{2q+1}+R,
 \qquad
 0\le R\le
 \frac{2z^{49}}{49(1-z^2)}
 \le\frac{9}{196\,3^{49}}<2^{-79}<Q^{-1}.
\]
Indeed, the expansion follows by integrating the geometric series for
$2/(1-z^2)$ from $0$ to $z$; bounding every remaining denominator by $49$
gives the displayed remainder bound.  The final strict numerical inequality
is the integer inequality $9\,2^{79}<196\,3^{49}$.
The code evaluates the first 24 terms by interval operations and adds one
integer unit to the upper endpoint.  This includes the omitted nonnegative
remainder.  The same procedure at $r=2$ encloses $\log2$.
For a positive interval, monotonicity of the logarithm allows its lower and
upper endpoints to be evaluated separately.  Boxing a positive rational
before this step therefore also gives a valid enclosure of its logarithm.

\paragraph{Exponentials and powers.}
To evaluate $e^x$ at an endpoint $x\le0$, choose $k\ge0$ so that
$z=-x/2^k\in[0,1/2]$.  The positive Taylor polynomial and its remainder obey
\[
 e^z=\sum_{n=0}^{20}\frac{z^n}{n!}+R,
 \qquad
 0\le R\le\frac{e^{1/2}(1/2)^{21}}{21!}<2^{-84}<Q^{-1}.
\]
Here $e^{1/2}<2$, for example by termwise comparison of its exponential
series with the geometric series, and $21!>2^{64}$; these two inequalities
give the strict $2^{-84}$ bound.  The verifier adds one integer unit to the
upper endpoint of the polynomial enclosure, takes its reciprocal, and then
squares $k$ times.  Each step is an enclosing interval operation, so the
result contains $e^x$.  Monotonicity supplies the enclosure for an interval
of nonpositive exponents.  For $0<\rho\le1$ and $0\le c\le1$, the expression
$\rho^c$ is evaluated as $\exp(c\log\rho)$.  The exact cases $\rho=1$, $c=0$,
and $\rho=0<c$ are handled directly.  The survival coefficient
$\rho^{1+c}$ is computed as $\rho\,\rho^c$.

\paragraph{Choice of tangent points and weighted planes.}
In the proof of Lemma~4.7, let $\phi_e$ denote the restriction of
$\mathcal L(\Delta,x,y)-\sigma x-\nu y$ to one of the three edges of the triangle.
For every chosen rational $u_e\in[0,1]$, convexity gives
\[
 \min_{0\le u\le1}\phi_e(u)
 \ge \phi_e(u_e)-u_e\phi'_e(u_e)+\min\{0,\phi'_e(u_e)\}.
\]
The verifier encloses the right-hand side and uses its lower endpoint.
Taking the minimum of the three resulting lower endpoints gives a certified
intercept $\beta$.  A floating-point derivative search is used only to
suggest the rational point $u_e$ (with denominator $2^{30}$).  The displayed
inequality holds for every point in $[0,1]$; neither root-finding accuracy nor
a floating-point assertion about an optimum is a premise of the proof.

If a rational slope is too large to be boxed at the chosen scale, the code
instead certifies the already weighted expression.  For the exact rational
weight $\omega\ge0$, it encloses $\omega\sigma$ and $\omega\nu$ directly and
applies the same tangent formula to $\omega\phi_e$.
Multiplication by a nonnegative constant preserves convexity, and the
three-edge argument in Lemma~4.7 remains valid.  Thus the resulting plane
has the form
\[
 \omega \mathcal L(\Delta,x,y)
 \ge \beta_{\mathrm w}+\omega\sigma x+\omega\nu y.
\]
Here $\beta_{\mathrm w}$ is the certified rational weighted intercept.
When $\omega>0$, dividing by $\omega$ gives the fixed valid unweighted
intercept $\beta=\beta_{\mathrm w}/\omega$; a zero-weight contribution is
discarded. Thus this computation also specifies a fixed plane that can be
multiplied by a variable proof weight in the manuscript's linear LP.
This is a scaled application of Lemma~4.7, rather than an additional
nonlinear lower-bound family.  Enclosing the exact weight at each arithmetic
operation, even when occurrences are treated independently, preserves
containment of the exact weighted formula.

\paragraph{Integer ranges for the supplied certificate.}
The range claim concerns the finite input files in this release.  It does
not assert that arbitrary malformed or adversarial input is accepted
safely.  Their grids have $N_p=10240$ and owner counts in
$\{160,320,640\}$, and each proposal has at most $1972$ planes.  All interval
endpoints are checked against $G$.  Products of two such endpoints have
magnitude less than $2^{124}$, and an endpoint multiplied by $Q$ has
magnitude less than $2^{118}$; both fit in a signed 128-bit integer.
Rational-to-interval scaling additionally has an explicit overflow check.
The rational routines cancel common factors and reject results outside
their 64-bit range.  Products of two 64-bit numerators or denominators,
and the two-product sums used in rational addition, fit in signed 128-bit
arithmetic for these inputs.

The few integer accumulators outside the interval class also remain in
range.  For example, there are at most $1972\cdot64$ range updates per
rectangle, with each update endpoint bounded in magnitude by $G$; even
counting both endpoint events, summing their absolute values, and then
multiplying by $N_p$ gives a bound less than $2^{94}$.  Suffix sums are
bounded by $N_pG<2^{76}$.  The exact common
parameter numerators are at most $20808318421$ with denominator $10^{10}$,
so the baseline and correction-prefix accumulations also fit easily in
128 bits.  In logarithm range reduction, the dyadic endpoints are in
$(0,G)$, and the reduced numerator plus denominator is less than
$96Q<2^{63}$.  The exponential range reduction requires a divisor of at
most $128$ under the same guard.  The completed replay checks every
rectangle without a range exception.  These finite range facts, together
with the directed formulas above, justify the arithmetic used for this
certificate.

\paragraph{Final comparison.}
The assembled coefficient intervals enclose a valid affine lower bound
from Lemma~B.1.  Replacing each coefficient by its lower endpoint preserves
a lower bound because every interval average lies in $[0,1]$.
Lemma~B.2 then reduces its minimum to the minimum of all suffix sums,
including the empty suffix.  If the resulting lower integer is $v$, the
target test is the exact integer comparison
\[
 2500v\ge1729Q.
\]
No rounded decimal, numerical optimization status, or floating-point
comparison is used to accept the competitive-ratio guarantee.

\section{Reproduction and the finite verification result}\label{sec:reproduce}
From the directory containing \texttt{replay.py}, run
\begin{verbatim}
python3 replay.py --audit-only
python3 replay.py --jobs 8
\end{verbatim}
The requirements are Python 3 and a GNU-compatible C++17 compiler supporting
\texttt{\_\_int128}. No numerical optimizer or external Python package is
needed. The first command checks the supplied file hashes, common parameters,
plane-mass budgets, assignments and complete coverage. The second also
compiles the verifier and reconstructs all 5120 directed arithmetic records.
Acceptance uses the integer comparison with $1729/2500$, not a displayed
floating-point decimal. Every reconstructed arithmetic field is compared
with the corresponding saved record.

The minimum reconstructed numerator at denominator $2^{56}$ is
$49839735470020217$. It occurs on
\[
 T_\ell=[10222/10240,10224/10240],
\]
with minimizing suffix index 20. Its exact margin above the target is
\[
 \frac{49839735470020217}{72057594037927936}-\frac{1729}{2500}
 =\frac{2939645868285289}{45035996273704960000}>0.
\]
The complete replay passes, and every time interval has lower bound at
least this minimum. Combining \eqref{eq:affine}--\eqref{eq:suffix} with this
finite computation proves \eqref{eq:certified}. Positivity and monotonicity
of the supplied drop rates give the asserted properties of $\theta$.
Finally, Lemma 4.6 of the manuscript transfers the uniform bound on
$\Gamma_{p^*}(\tau)$ to the expected reward of the algorithm.

The deterministic final assembly can be repeated separately:
\begin{verbatim}
python3 generation/rebuild_proposals.py \
  --output-dir regenerated_proposals --compare-dir proposals
\end{verbatim}
It reconstructs all 5120 final files from 165 frozen candidate files and
646 seed intervals, checking the output hashes and, with the displayed
option, byte equality. This reconstructs the supplied final certificate
from its frozen inputs; it does not reproduce the earlier exploratory
optimization that discovered those inputs. Verification of feasibility
does not depend on that discovery process.

\paragraph{Files implementing the proof.}
\texttt{sources/strict160\_verify.cpp} constructs the complete-rectangle
coefficients and enumerates the suffixes;
\texttt{sources/log\_plane.hpp} certifies the three-edge tangent intercepts;
\texttt{sources/interval.hpp} supplies the directed arithmetic.
\texttt{parameters.json}, \texttt{proposals/}, \texttt{intervals.json} and
\texttt{records.json} specify the candidate and the full finite verification
records. \texttt{certificate.json} records the integrity hashes.
The legacy \texttt{tail} field in proposal files is parsed but is not used
to determine the route; the route is exactly the one specified in
Section~\ref{sec:allocation} and in the verifier's \texttt{route()} function.

\paragraph{Compiling this supplement.}
The PDF is generated from this single source using
\begin{verbatim}
pdflatex -interaction=nonstopmode -halt-on-error Verification.tex
pdflatex -interaction=nonstopmode -halt-on-error Verification.tex
\end{verbatim}
Only standard LaTeX packages are used. No external bibliography, figure,
input file or internet access is needed.
\end{document}
